\section{Limitations of Conventional Front-Side Power Delivery Networks}

Conventional integrated circuits rely on front-side power delivery (FSPDNs) where power and signal interconnects are routed via the same back-end-of-line (BEOL) metal stack. In this layout, supply voltages are distributed from the package connections through all individual metal layers, down to the transistor level. This architecture has been used for multiple generations of integrated circuits, however continued scaling below 5 nm has revealed several limitations that have made reliable power delivery more difficult \cite{9817394},\cite{10244504}. 

As a result of the scaling of interconnects, an issue arises from the increased resistance. As metal dimensions shrink, resistivity increases due to surface and grain-boundary scattering, leading to larger IR~drops across the power grid\cite{8993617}. At the local metal layers closest to the transistors, wire widths fall below 20\,nm, and at these dimensions copper resistivity rises sharply because the wire is narrower than the electron mean free path\cite{8932458}. This results in a reduced effective supply voltage at the transistors, which can degrade circuit performance and reliability. In the most advanced node, highly resistive local metal layers contribute significantly to IR-drop hotspots, particularly in areas with dense logic \cite{8993617}. Studies conducted with an ARM Cortex-A53 design at a 3\,nm technology node show that even with an optimised power grid, conventional FSPDNs did not meet the 10\% V\textsubscript{DD} IR-drop target\cite{9817394}.

Another key limitation that arises due to scaling is the increased competition for the limited routing resources available for power and signal interconnects. In FSPDNs, both power rails and signal interconnects occupy the same metal stack. As such, designers must dedicate a large portion of routing resources to maintain a sufficiently reliable power grid that is able to provide the necessary performance. This results in reduced routing resources for signal interconnects and increased routing congestion. In high-performance designs, the power grid can take up four or more metal layers that would otherwise carry signals, which forces designers to choose between a larger die area or lower performance\cite{10631463}. Block-level studies at imec showed that high-performance standard-cell libraries need an 18-CPP power stripe pitch on the front side just to keep IR~drop acceptable, while backside delivery reaches the same power integrity without using any front-side routing tracks at all\cite{10631463}.

To mitigate these limitations, Buried Power Rails (BPR) are used as a current solution. These designs move power rails beneath active devices in order to reduce resistance and make more routing resources available to the signal interconnect. BPR-based designs have shown a 25\% reduction in on-chip IR drop and 17\% reduction in off-chip voltage droop compared with conventional front-side PDNs \cite{9817394}. However, even with these improvements, standalone BPR designs are viewed as intermediate solutions\cite{10244504}. The fundamental limitations of front-side routing remain.

Backside power delivery networks (BSPDNs) have been explored as a long-term solution to these issues. In this architecture, the entire power grid is moved to the backside of the wafer, typically through nano through-silicon vias (nTSVs) that connect the backside metal layers directly to the transistor contacts. By fully separating power and signal routing onto opposite sides of the die, BSPDN architecture eliminates the resource conflict between the two and enables both lower IR~drop and improved routing density simultaneously. Multiple studies have confirmed the viability of this approach at 3\,nm and below\cite{10631463},\cite{8932458},\cite{10244504}.

